\section{One-Hop Scheduling and Implementation Details}
\label{app:one-hop}
Section~\ref{sec:one-hop} states the schedule and its observable contract.
This appendix preserves the dispatch algorithm, proof, worked quotas, input
lifetime argument, and module-specific processing details.

\subsection{Balanced dispatch and worked quotas}
Implement Equation~\eqref{eq:balanced} with a quotient/remainder accumulator.
Quotient and
remainder $a=\floor{W/H}$ and $r=W\bmod H$ are refreshed at configuration,
reset, and the transition from a rebuilt window to a clean cache.
Starting the accumulator $e$ at zero, each sample executes $a$ units plus one
when $e+r\ge H$, then retains $(e+r)\bmod H$:
\begin{lstlisting}[float=htbp,language=C++,caption={Balanced one-hop dispatch. execute\_next\_unit retains stage and traversal state; output stays private until publication.},label={lst:whole-hop}]
// At a frame boundary: latch settings; compute a = W/H, r = W%H.
// Each active sample call:
capture_input();
if (phase == 0) { latch_input_endpoint(); e = 0; begin_frame(); }
size_t quota = a;
e += r;
if (e >= H) { e -= H; ++quota; }
for (size_t i = 0; i < quota; ++i) execute_next_unit();
if (++phase == H) {
    publish_owned_output(); phase = 0;
    if (window_was_dirty) use_clean_frame_quota();
}
\end{lstlisting}

\begin{proposition}[Exact completion and per-call work]
For $W>0$ and integer $H\ge1$, the ordered analysis completes on the $H$th
call, with at most $\ceil{W/H}$ units on any call. For $H=\rho N$ with fixed
$\rho>0$, this maximum is $O(\log N)$.
\end{proposition}
\begin{proof}
After $s$ calls the accumulator rule has issued
$sa+\floor{sr/H}=\floor{sW/H}$ units. At $s=H$ this is $W$;
at $s=H-1$ it is strictly smaller. Each difference is either $a$ or $a+1$.
Substitution of Equation~\eqref{eq:wholework} gives the asymptotic bound.
\end{proof}
For a clean $N=2048,H=1024,o=1$ frame, $W=8194$ and the maximum quota is nine:
1022 calls execute eight units and two execute nine. With a dirty window,
$W=11266$: 1022 calls receive eleven units and two receive twelve.
At most $\ceil{\ceil{W/H}/4}$ preparation pairs execute in a dirty call,
whereas subsequent stages may execute up to $\ceil{W/H}$ operations.
For Fourier ($o=2$), the corresponding counts are $W=9219$ for a clean cache
and $W=12291$ for a dirty cache. Output calls execute at most
$\ceil{\ceil{W/H}/o}$ bins, limiting the additional concentration of costly
coordinate mapping during a window rebuild. Reweighting redistributes bursts
between stages; it does not reduce the transform arithmetic. Dense dispatch skips
credit over a whole segment, and SIMD butterflies share traversal setup
within the same quota. The bound includes window/cache
rebuilds, but excludes construction and host lifecycle callbacks. Constant
input/control work and frame configuration remain, including an $O(\log N)$
plan-index calculation. No $O(N)$ boundary pass remains in steady-state engine
analysis. This count assumes fixed-width scalar/SIMD arithmetic and the bounded
module output callbacks; the generic API cannot constrain an arbitrary callback.
It does not assert equal duration across units or bound OS interruptions.

The current implementation also omits multiplication by the exactly unity
first twiddle in dense SIMD butterfly segments. This reduces arithmetic
without changing scheduling credit or the dependency order. It is a standard
FFT simplification, not a new transform algorithm.

\subsection{Retained-input proof and plan lifecycle}
Incremental packing cannot read a ring that has already overwritten its
selected samples. With prepared maxima $N_{\max},H_{\max}$, the production
ring has capacity $C=N_{\max}+H_{\max}$. At selection, the oldest needed sample
is $N-1$ samples behind the newest. At most $H-1$ new captures occur before
completion. Its maximum lag is therefore $N+H-2<C$, so it cannot be
overwritten before its packing unit. This conservative bound remains valid
when capture pauses. The ring and all transform/output arrays are preallocated;
no whole-frame capture copy is required.

A maximum-size bit-reversal table serves shorter transforms by right shifts;
a maximum-size scalar twiddle table serves them by strides (Appendix~\ref{app:implementation}). SIMD lanes share
those scalar twiddles. Window weights and smoothing endpoints rebuild only in
the scheduled units that use them. Length changes logically clear input and
EMA history. Reset cancels partial work and marks partially rebuilt caches
invalid. Increasing the supported hop during a sample-rate callback can grow
the ring and allocate; this is outside the allocation-free sample path.

\subsection{Smoothing and output ownership}
When frequency smoothing is enabled, for magnitudes $a_k=|X[k]|$,
reconstruction builds $P_0=0$ and
$P_{k+1}=P_k+a_k$. Inclusive smoothing bounds $[\ell_k,h_k]$ then give
\begin{equation}
 v_k=\frac{P_{h_k+1}-P_{\ell_k}}{h_k-\ell_k+1},\qquad
 y_r[k]=\alpha y_{r-1}[k]+(1-\alpha)v_k.
\end{equation}
The octave interval is defined in Appendix~\ref{sec:cost}: endpoints are floored
to bins; if its upper frequency exceeds Nyquist, it is set to Nyquist and the
lower endpoint is set to that upper frequency divided by $2^\beta$. With
frequency smoothing off, $v_k=a_k$ and prefix/band work is skipped. With
$\alpha=0$, the implementation stores $y_r[k]=v_k$ directly, retaining that
history for later averaging. When enabled, the full prefix stage precedes output
because a low bin's interval may require higher bins. These are magnitude
means, not power means or constant-Q analysis.

For Average control $T>0$ seconds, the retained convention is
$\alpha=\exp(-10\Delta t/T)$; $T=0$ disables averaging. Spectre uses $\Delta
t=H/f_s$. The Fourier module uses its Hop control in seconds and truncates
$\Delta t f_s$ to $H$, retaining a sub-sample quantization difference between
the control and actual interval. The previous $L/H$ completion bias is removed.
Window coherent-gain factors and module-specific display scaling are preserved;
this scheduling change does not redefine amplitude units.

Fourier additionally caches the frequency coordinate and slope gain for each
bin in engine-owned storage. A geometry change invalidates a prefix counter
in constant time; each scheduled output computes at most one missing entry.
Only consecutively initialized entries advance that counter, so cancellation
cannot expose stale entries. Magnitude scaling still evaluates every current
value. This uses approximately 64\,KiB at $N_{\max}=16384$, with unchanged
frame timing and no approximation of the coordinate formulas. Reuse reduces
steady-state work; changing geometry every frame removes that reuse. These
engineering changes postdate the report's historical timing campaigns and
must be measured separately rather than attributed to those data.

Each output unit maps a four-channel bin into the Fourier module's curve
points, or writes one Spectre column magnitude. Only after all units finish
does the engine publish the destination. A single-producer/single-consumer
triple-slot mailbox uses acquire/release atomic exchange to transfer ownership
of a completed slot and obtain a reusable slot. The consumer holds its slot
until its next successful exchange; the producer never overwrites that slot. A
slow consumer may miss intermediate snapshots. The Fourier module has one curve
mailbox; the Spectre module has one mailbox per history column. Spectre's full
image history and recoloring belong to the UI, outside the engine schedule.

Freeze semantics remain distinct: the Fourier module stops capture but
continues analysis of retained input, whereas Spectre pauses both acquisition
and the schedule and resumes its partial hop. Continuous-input timestamp
equations do not apply across such pauses. Host sample-rate changes cancel
partial analysis; Spectre reset still clears and publishes its history columns.
These lifecycle costs, UI work, other shared controls, and host scheduling
prevent a broad claim of hard real-time safety.
